\documentclass[12pt]{article}
\usepackage[utf8]{inputenc}
\usepackage[T1]{fontenc}
\usepackage{hyperref}
\title{Zemgalian Language Resources for K-12 Curriculum Development}
\author{Compiled by Claude Sonnet 5.5 (claude-sonnet-5-5) for Andr\'as Kornai}
\date{}
\begin{document}
\maketitle

\section{Introduction}
Zemgalian (Semigallian, ISO 639-3 xzm) is an extinct East Baltic language, thought to have died out by about the 16th century. The only evidence is place names, personal names and possible substrate influence on Latvian and Lithuanian dialects. The search found no revival movement, dictionary, course or school materials. Only reference and scholarly pages turned up.

\section{Primary Educational Resources}
None were found.

\section{Linguistic and Reference Materials}
\begin{itemize}
\item Encyclopedic overview: \url{https://en.wikipedia.org/wiki/Semigallian_language}
\item Background on the region: \url{https://en.wikipedia.org/wiki/Semigallia}
\item Lithuanian article ``\v Ziemgali\u{} kalba ir jos reliktai'' on the language and its relics: \url{https://lituanistika.lt/content/30803}
\item Lithuanian article on the language's legacy in Lithuanian Zemgale (\emph{\v Ziemgala}, 2021): \url{https://lituanistika.lt/content/109050}
\end{itemize}

\section{Religious and Community Resources}
None were found.

\section{Digital and Audio Resources}
None were found. A Zemgale tourism guide has an entry on the Semigallian language, but only its table of contents was visible: \url{https://zemgale.lv/lv/media/1335/download}

\section{Curriculum Development Context}
There is almost no linguistic material to teach from. A K-12 unit could cover Semigallian history, place names and Baltic language history using the sources above, and would need a Baltic linguist's input. Searches in Latvian or Lithuanian would be the next step.
\end{document}